\documentclass[11pt]{article}

\title{Building a GPU Job Platform from Lab PCs\\
A Technical Guide to distributetr9}
\author{Written alongside the code, November 2026}
\date{}

\begin{document}

\maketitle

\begin{abstract}
This guide explains how distributetr9 works and why it is built the way it
is. distributetr9 turns a room of ordinary Windows PCs with consumer GPUs into
a shared cluster: students submit a Python training script in a web portal, an
idle lab PC runs it inside a Docker container, and the logs and outputs come
back to the student. Along the way the project touches most of the ideas behind
real distributed systems: leases, fencing tokens, idempotency, database
locking, containers, image registries, signed software updates, HTTPS, chaos
testing and operations. Each chapter introduces one idea with a small example
you can run or picture, shows where it lives in this repository, and lists
other places the same idea shows up in industry. The final chapters retell the
real bugs found while building the project, because debugging stories teach
more than finished code does.
\end{abstract}

\section{The Big Picture}

\subsection{The problem}

A university lab has dozens of Windows desktops with NVIDIA graphics cards.
Most of the day they sit idle. Meanwhile students training neural networks for
coursework or a thesis have no GPU at all. distributetr9 connects the two:

\begin{enumerate}
\item A student uploads \texttt{train.py} (plus optional \texttt{requirements.txt})
      in a web portal and asks for, say, ``a GPU with at least 6 GB''.
\item A central server, the \textbf{control plane}, stores the job and puts it
      in a queue.
\item A small program on every lab PC, the \textbf{agent}, keeps asking the
      server ``do you have work for me?''. A free PC with a matching GPU gets
      the job.
\item The agent runs the script inside a locked-down \textbf{Docker
      container}, streams the output lines back live, and uploads the files
      the script wrote.
\item The student watches the logs in the browser and downloads the results.
\end{enumerate}

That sounds simple. The hard part, and the reason most of the code exists, is
what happens when things go wrong: a PC reboots in the middle of a job, Wi-Fi
drops for a minute, a student cancels at the exact moment the job finishes, or
two PCs both think they own the same job. The project's specification puts it
in one sentence: \emph{the hard problem is not the portal, it is keeping exactly
one accepted result when nodes disconnect, leases expire, completions are
retried, and cancellation races completion.}

\subsection{The components}

\begin{verbatim}
   Student's browser
        |  HTTPS
        v
  +-----------+      +---------------------------------------------+
  |  Portal   | ---> |              Control plane (FastAPI)        |
  | (Next.js) |      |  REST API  |  scheduler  |  lease worker    |
  +-----------+      +------+-------------+-----------------+------+
                            |             |                 |
                      +-----v----+  +-----v-----+    +------v------+
                      | Postgres |  |   MinIO   |    |  Registry   |
                      | (state)  |  | (files,   |    | (container  |
                      |          |  |  logs)    |    |  images)    |
                      +----------+  +-----------+    +------^------+
                            ^                               |
              HTTP: register, heartbeat, claim,             | docker pull
              renew, logs, outputs, complete                |
                            |                               |
   +------------------------+----------+--------------------+------+
   |  Lab PC 1 (Windows + WSL Ubuntu)  |  Lab PC 2  ...  Lab PC 60 |
   |  launcher -> agent -> docker run  |                           |
   +-----------------------------------+---------------------------+
\end{verbatim}

\begin{description}
\item[Portal] A Next.js web app. It never talks to the database; it forwards
  every request to the control plane's API.
\item[Control plane] A Python FastAPI service. It owns all state and makes
  every decision: who may log in, which job goes to which PC, whether a result
  is accepted.
\item[PostgreSQL] The source of truth for users, jobs, attempts and nodes.
  Its row locks are what make the scheduler safe when many PCs ask for work at
  once.
\item[MinIO] An S3-compatible object store for submitted archives, log
  chunks and output files.
\item[Registry] A Docker image registry (\texttt{registry:2}). Custom Python
  environments are built once and pulled from here by every PC.
\item[Agent] A Python program with no third-party dependencies that runs in
  WSL Ubuntu on each lab PC. It registers, heartbeats, claims work and runs it.
\item[Launcher and supervisor] A Windows scheduled task starts WSL at boot
  without anyone logging in; a small bash launcher keeps the agent running and
  rolls back bad updates.
\end{description}

\subsection{Two kinds of node: trainers and builders}

Every lab PC runs the same agent program, in one of two roles chosen at install
time:

\begin{description}
\item[Training nodes] (\texttt{role: train}, the default) run students' jobs.
  They never build environments.
\item[Build nodes] (\texttt{role: build}) build custom environments from
  \texttt{requirements.txt}, test them and push them to the registry. They never
  run training jobs.
\end{description}

The agent reports its role when it registers, and the same claim endpoint
hands out different work depending on it:

\begin{lstlisting}[language=Python]
# control/app/routers/agent.py (simplified)
@router.post("/agent/claim")
async def claim_job(node: Node = Depends(get_current_node), ...):
    if is_train_node(node):
        attempt = await claim(session, node)          # a training job
        ...
    else:
        build = await claim_build(session, node)      # an environment build
        ...
\end{lstlisting}

Why separate them? A build downloads packages from the internet, compiles
wheels and pushes multi-gigabyte images; a training job wants a quiet machine
with its whole GPU. Keeping them apart means a slow \texttt{pip install} never
delays someone's training run, and only the builder needs internet access and
registry push credentials. In this repository the role is chosen with
\verb|--role build| on the agent or installer, \verb|-Role build| on the
Windows bootstrap, or \verb|LAB_NODE_ROLE=build| in the environment.

\paragraph{A current limitation.} A builder only takes builds for its own GPU
class, or builds marked \texttt{cpu}. Jobs that ask for a GPU class and bring a
\texttt{requirements.txt} produce a build for that class, so a CPU-only builder
cannot serve them; they wait in \texttt{building} until a builder with that GPU
class exists. Catalog environments need no builder at all.

\subsection{Life of a job}

Following one job end to end is the fastest way to understand the system.

\begin{enumerate}
\item \textbf{Submit.} The portal sends \texttt{POST /jobs} with the archive.
  The control plane checks the archive, stores it in MinIO under
  \texttt{submissions/<job-id>/input.tar.gz}, and inserts a \texttt{jobs} row
  with \texttt{desired\_state = run}.
\item \textbf{Claim.} Every agent calls \texttt{POST /agent/claim}. The call is
  a \emph{long poll}: the server holds it open for up to 25 seconds until a job
  fits. When one fits, the server inserts a \texttt{job\_attempts} row with a
  fresh random \textbf{fencing token} and a \textbf{lease} that expires in 30
  seconds.
\item \textbf{Run.} The agent downloads the input, starts \texttt{docker run},
  and every few seconds calls \texttt{/renew} to extend the lease. Log lines are
  posted in numbered chunks.
\item \textbf{Upload.} Output files are uploaded under an
  \emph{attempt-specific} prefix, not the final location.
\item \textbf{Complete.} The agent calls \texttt{/complete} with its fencing
  token. The server checks the token is still the current one, verifies every
  output checksum, copies the files to \texttt{jobs/<id>/promoted/}, and marks
  this attempt as \emph{the} accepted result.
\item \textbf{Failure path.} If the PC vanishes, the lease runs out. A
  background worker marks the attempt \texttt{unknown}, then \texttt{failed}
  with category \texttt{lease\_lost}, and the job goes back in the queue for
  another PC. If the old PC comes back and tries to complete, its token no
  longer matches and the server answers \texttt{409 fenced}.
\end{enumerate}

\subsection{Words used throughout}

\begin{description}
\item[Node] A lab PC running the agent.
\item[Job] What the student asked for. A job can have several attempts.
\item[Attempt] One try at running a job on one node. At most one attempt per
  job is ever \emph{accepted}.
\item[Lease] A time-limited claim on a job. Renew it or lose it.
\item[Fencing token] A random value that identifies the current attempt; stale
  attempts are rejected because they hold an old token.
\item[Compatibility class] A label like \texttt{ampere-8g} or
  \texttt{turing-6g} describing a GPU's architecture and memory band.
\item[Environment] A container image a job runs in: a maintained catalog image
  or one built from the student's \texttt{requirements.txt}.
\end{description}

\section{How the Repository Is Organised}

\begin{verbatim}
distributetr9/
  control/                FastAPI control plane
    app/
      main.py             app startup, background workers
      config.py           every setting, read from environment variables
      models.py           SQLAlchemy tables
      db.py               engine, sessions, migrations on startup
      scheduler.py        claim, renew, complete, lease expiry
      envbuild.py         custom environment builds
      agent_bundle.py     signed agent updates, install package
      alerts.py ops.py    operations: alerts, garbage collection
      routers/            HTTP endpoints: portal.py, agent.py, auth.py
    migrations/           Alembic schema migrations
  agent/                  runs on each lab PC (Python standard library only)
    agent.py              main loop: register, heartbeat, claim, run
    executor.py           docker run / local python, log pumping
    classify.py           GPU detection and class mapping
    updater.py            verify and install signed updates
    launcher.sh run.sh    process supervision and rollback
    classes.json          GPU class catalog
  portal/                 Next.js web UI
  harness/                fake agents for tests and load testing
  tests/                  pytest suite (runs against real Postgres + MinIO)
  sim/                    simulated five-PC lab + failure exam
  deploy/                 production compose, Caddy, backup, restore
  scripts/                dev launcher, lab PC installer, smoke test
  supervisor/             Windows boot task and one-command bootstrap
  spec/                   the design specification (14 modules, 12 phases)
  docs/                   runbooks, registry notes, this guide
\end{verbatim}

\subsection{Why this split}

Each directory matches one \emph{deployable thing} or one \emph{audience}:

\begin{itemize}
\item \texttt{control/} and \texttt{portal/} run on the server.
\item \texttt{agent/}, \texttt{supervisor/} and the installer run on lab PCs.
  The agent deliberately uses only the Python standard library, so installing
  it means copying files, and updating it means copying newer files. There is
  no \texttt{pip install} on sixty machines that might each fail differently.
\item \texttt{harness/}, \texttt{tests/} and \texttt{sim/} exist only to
  prove the other parts work.
\item \texttt{spec/} was written before the code. It froze the scope of version
  1 and split the work into phases, each ending with a \emph{gate}: a concrete
  demonstration that must pass before the next phase starts.
\end{itemize}

\subsection{Configuration: one settings object}

Every tunable value lives in \texttt{control/app/config.py} as a field on a
\texttt{pydantic-settings} class. Pydantic reads each field from an environment
variable of the same name (case-insensitive), falling back to a
\texttt{.env} file and then to the default:

\begin{lstlisting}[language=Python]
class Settings(BaseSettings):
    model_config = SettingsConfigDict(env_file=".env", extra="ignore")

    lease_ttl_seconds: int = 30
    claim_poll_seconds: int = 25
    agent_rollout_ring: int = 2

settings = Settings()   # LEASE_TTL_SECONDS=10 in the environment wins
\end{lstlisting}

The benefit is that tests, the simulated lab and production run the same code
with different numbers. The simulated lab, for example, sets a 10 second lease
so that failure scenarios resolve in seconds rather than minutes.

\paragraph{Other use cases.} This is the ``config in the environment'' rule of
the Twelve-Factor App methodology, used by most containerised services. The
same pattern appears in Django (\texttt{django-environ}), Spring Boot
(\texttt{application.properties} overridden by environment variables) and
Kubernetes ConfigMaps.

\subsection{The development workflow}

\begin{lstlisting}[language=bash]
make dev        # Postgres, MinIO, registry, API :8080, portal :3000, agents
make test       # pytest against a throwaway lab_test database
make lint       # ruff: syntax errors, undefined names, unused imports
make sim        # five simulated lab PCs in Docker
make sim-exam   # break them in thirteen different ways
\end{lstlisting}

A \texttt{Makefile} here is only a list of named shortcuts; it does not build
anything. It helps newcomers find the important commands without reading the
README.

\section{Distributed Systems Ideas}

This chapter covers the ideas that make the platform correct when machines and
networks fail. They are worth learning well: they appear in databases, job
queues, payment systems and cloud schedulers.

\subsection{Pull, not push}

There are two ways to hand out work. In a \emph{push} design the server picks
a worker and sends it the job. In a \emph{pull} design each worker asks for
work when it is ready. distributetr9 pulls:

\begin{lstlisting}[language=Python]
# agent/agent.py (simplified)
while True:
    status, hb = self.heartbeat()          # "I am alive, here is my state"
    status, payload = self.claim()         # "Any work for me?" (long poll)
    attempt = payload.get("attempt")
    if attempt:
        self.execute(attempt)
\end{lstlisting}

Pulling suits lab PCs well:
\begin{itemize}
\item The server never needs to open a connection \emph{to} a PC. PCs sit
  behind Windows firewalls and NAT; outgoing HTTP from the PC almost always
  works.
\item A busy, broken or rebooting PC simply stops asking. The server does not
  need to track whether a push succeeded.
\item The PC decides when it is ready, for example not while a student is
  gaming on it (the \texttt{gpu\_desktop\_busy} check).
\end{itemize}

\paragraph{Long polling.} Asking every second wastes requests; asking every
minute adds a minute of delay. A \emph{long poll} gets both: the agent asks,
and the server holds the request open, checking once a second, until work
appears or 25 seconds pass. The response comes back the moment a job fits.

\paragraph{Other use cases.} GitHub Actions self-hosted runners, GitLab
runners, Jenkins agents in ``inbound'' mode, Celery and Sidekiq workers pulling
from a queue, and Amazon SQS long polling (\texttt{WaitTimeSeconds}) all pull.

\subsection{Heartbeats}

Every few seconds the agent sends \texttt{POST /agent/heartbeat} with its free
disk space, whether the desktop GPU is busy, and whether it is healthy. The
server records \texttt{last\_heartbeat\_at}. A background worker marks nodes
whose heartbeat is older than three intervals as \texttt{offline}, and the
scheduler skips them. The heartbeat \emph{response} carries commands back:
\texttt{drain}, \texttt{quarantine}, or an available agent update. That gives
the server a way to talk to PCs without ever connecting to them.

\subsection{Leases}

A lease is ownership with an expiry date. When a node claims a job it gets the
attempt for 30 seconds. To keep it, the node must renew before time runs out.

\begin{verbatim}
 t=0s    claim       lease_expires_at = 30s
 t=10s   renew       lease_expires_at = 40s
 t=20s   renew       lease_expires_at = 50s
 t=25s   PC loses power ...
 t=50s   lease expires         -> attempt state: unknown
 t=110s  + disconnect grace    -> failed (lease_lost), job re-queued
\end{verbatim}

Why not just wait for the PC to say it failed? Because a crashed PC cannot say
anything. A lease turns ``silence'' into a decision the server can make on its
own clock. The grace period avoids giving up on a PC that only missed a few
renewals during a Wi-Fi hiccup.

\paragraph{Other use cases.} DHCP leases your laptop an IP address; Kubernetes
leader election uses Lease objects; distributed locks in etcd, ZooKeeper and
Redis (Redlock) are leases; Amazon SQS's ``visibility timeout'' is a lease on a
message.

\subsection{Fencing tokens: rejecting zombies}

Leases create a new problem. Picture this:

\begin{enumerate}
\item PC A claims job 7 and starts training.
\item PC A's network drops for two minutes. Its lease expires.
\item The server gives job 7 to PC B, which finishes and completes it.
\item PC A's network returns. It finishes too and calls \texttt{/complete}.
\end{enumerate}

PC A is a \emph{zombie}: it believes it owns a job it lost long ago. If the
server accepted its result, job 7 would have two results, or the wrong one.

The fix is a \textbf{fencing token}. Each attempt gets a fresh random UUID when
it is created. Every call that changes state must carry it, and the server
compares it with the token stored on the attempt:

\begin{lstlisting}[language=Python]
# control/app/scheduler.py
def _require_token(attempt: JobAttempt, token: uuid.UUID) -> None:
    if attempt.fencing_token != token:
        raise AppError(409, "fenced", "stale fencing token", ...)
\end{lstlisting}

PC A's attempt was failed when its lease expired, and job 7's accepted attempt
is now PC B's. PC A's \texttt{/complete} hits the ``attempt superseded'' check
and gets \texttt{409}. The zombie's work is discarded, never promoted. The
simulated failure exam checks exactly this in its \texttt{network\_drop}
scenario.

\paragraph{Other use cases.} Martin Kleppmann's essay ``How to do distributed
locking'' popularised the term. Apache Kafka uses producer epochs, HDFS uses
generation stamps, and Raft and Paxos use terms and ballot numbers: all are
fencing tokens.

\subsection{Boot sessions}

A PC that reboots gets a brand new agent process, but the same node identity.
On registration the server issues a new \texttt{boot\_session\_id} and fails
every attempt the node still owned (\texttt{on\_node\_boot}). Requests carrying
an old boot session get \texttt{409 stale\_boot\_session}. This catches a case
leases alone would only catch slowly: a PC that rebooted in ten seconds has
certainly lost its running container, so there is no point waiting a full lease.

\subsection{Idempotency: safe to repeat}

Networks lose responses. If an agent sends \texttt{/complete} and the reply
never arrives, it cannot know whether the server processed it. The only safe
move is to send it again, so every state-changing call must be
\textbf{idempotent}: doing it twice has the same effect as doing it once.

\begin{itemize}
\item \texttt{/complete} twice with the same token and exit code returns the
  stored result. With a \emph{different} exit code it returns \texttt{409},
  because that signals a real bug.
\item \texttt{/claim} twice from the same node returns the node's existing
  attempt instead of a second one.
\item \texttt{POST /jobs} accepts an \texttt{Idempotency-Key} header. A browser
  that retries a submit after a timeout gets the same job back instead of a
  duplicate:
\end{itemize}

\begin{lstlisting}[language=Python]
# control/app/routers/portal.py (simplified)
if idem_key:
    existing = await session.scalar(
        select(Job.id).where(Job.user_id == user.id,
                             Job.idempotency_key == idem_key))
    if existing is not None:
        return serialize_job(await load_job(session, existing))
\end{lstlisting}

A unique index on \texttt{(user\_id, idempotency\_key)} is the backstop: if two
identical requests race past the check at the same moment, the second insert
fails and the handler returns the job the first one created.

\paragraph{Other use cases.} Stripe's API is the textbook example: every
payment request takes an \texttt{Idempotency-Key} so a retried charge never
bills twice. HTTP itself defines \texttt{PUT} and \texttt{DELETE} as idempotent.

\subsection{At-least-once execution, exactly-once acceptance}

Truly running something \emph{exactly once} across unreliable machines is
impossible in general. What the platform guarantees instead is:

\begin{itemize}
\item A job may \emph{run} more than once (a zombie may keep training after it
  was replaced). That is \emph{at-least-once execution}.
\item Exactly one attempt's result is ever \emph{accepted and promoted}. That
  is \emph{exactly-once acceptance}, enforced by fencing and a single
  \texttt{accepted\_attempt\_id} column on the job.
\end{itemize}

Message queues make the same trade: Kafka and SQS deliver at least once, and
consumers make processing idempotent to get an exactly-once \emph{effect}.

\subsection{State machines}

An attempt moves through named states, and only some moves are legal:

\begin{verbatim}
 leased -> starting -> running -> uploading -> completed
    \          \          \          \
     +----------+----------+----------+--> failed | cancelling -> cancelled
                       any active state --> unknown (lease expired)
                       unknown --> back to running (renewed in time) | failed
\end{verbatim}

\begin{lstlisting}[language=Python]
# control/app/enums.py
FORWARD_TRANSITIONS = {
    "leased":    {"starting", "running", "uploading", "failed",
                  "cancelling", "unknown"},
    "running":   {"uploading", "failed", "cancelling", "unknown"},
    "uploading": {"completed", "failed", "cancelling", "unknown"},
    ...
}
\end{lstlisting}

Writing transitions down as data means an impossible report, such as
\texttt{completed} going back to \texttt{running}, becomes a clear
\texttt{409 illegal\_transition} instead of silently corrupting state.

\paragraph{Other use cases.} TCP connections, order fulfilment, CI pipelines,
Kubernetes pod phases and AWS Step Functions are all explicit state machines.

\section{The Database: Locks, Transactions and Migrations}

\subsection{Why PostgreSQL does the hard part}

Sixty PCs can ask for work in the same second. If two of them read ``job 7 is
free'' at the same time, both may take it. Instead of inventing a locking
scheme, the scheduler lets PostgreSQL do it with \textbf{transactions} and
\textbf{row locks}. Each request runs in one transaction: either all of its
changes are saved together, or none are.

\subsection{SELECT \dots{} FOR UPDATE SKIP LOCKED}

\texttt{SELECT ... FOR UPDATE} reads rows \emph{and locks them} until the
transaction ends. Another transaction trying to lock the same row waits.
Adding \texttt{SKIP LOCKED} changes ``wait'' into ``skip it and look at the
next row''. That one clause turns a table into a safe work queue:

\begin{lstlisting}[language=SQL]
-- Two workers run this at the same moment.
BEGIN;
SELECT id FROM jobs
 WHERE desired_state = 'run' AND accepted_attempt_id IS NULL
 ORDER BY priority DESC, created_at
 LIMIT 1
 FOR UPDATE SKIP LOCKED;     -- worker A locks job 7, worker B skips to job 8
-- ... create the attempt ...
COMMIT;                      -- lock released
\end{lstlisting}

\subsection{A real bug: locking the whole queue}

The first version of \texttt{claim()} locked \emph{every} waiting job, then
picked one in Python:

\begin{lstlisting}[language=Python]
# before: locks all queued jobs for the whole claim
stmt = (select(Job)
        .where(Job.desired_state == "run",
               Job.accepted_attempt_id.is_(None))
        .with_for_update(skip_locked=True))
jobs = (await session.scalars(stmt)).all()
\end{lstlisting}

With two PCs asking at the same moment, PC A locked all twenty queued jobs
while it decided. PC B skipped every locked row, saw an empty queue, and went
back to sleep, even though nineteen jobs were waiting. Nothing broke, so no
test noticed; the cluster just ran slower than it should.

The fix reads candidates \emph{without} locks, decides in Python which job
fits this node, and then locks \emph{only that one row}. If another node got it
first, it moves to the next candidate:

\begin{lstlisting}[language=Python]
# after (control/app/scheduler.py, simplified)
jobs = (await session.scalars(candidates_stmt.limit(200))).all()
for job in jobs:
    if not node_fits_job(node, job, envs.get(job.environment_id)):
        continue
    locked = await session.scalar(
        select(Job).where(Job.id == job.id,
                          Job.accepted_attempt_id.is_(None))
                   .with_for_update(skip_locked=True))
    if locked is None:          # another node won this one
        continue
    chosen = locked
    break
\end{lstlisting}

The test that proves it opens two database sessions, starts a claim in the
first \emph{without committing} (so its lock stays held), and checks the second
still gets the other job. On the old code the second claim returned nothing.

\subsection{Partial unique indexes}

``A node runs at most one job at a time'' is enforced by the database itself,
not only by Python:

\begin{lstlisting}[language=SQL]
CREATE UNIQUE INDEX uq_attempts_one_active_per_node
    ON job_attempts (node_id)
 WHERE observed_state IN ('leased','starting','running',
                          'uploading','cancelling','unknown');
\end{lstlisting}

A \emph{partial} index only covers rows matching its \texttt{WHERE}. Finished
attempts are not in it, so a node can have hundreds of past attempts but only
one active. If a bug ever tried to insert a second active attempt, Postgres
would reject the insert. The idempotency key uses the same trick with
\texttt{WHERE idempotency\_key IS NOT NULL}.

\paragraph{Other use cases.} ``One active subscription per user'', ``one
default address per customer'', ``one open shift per employee''.

\subsection{Talking to Postgres from Python: SQLAlchemy async}

The control plane uses SQLAlchemy 2 with the \texttt{asyncpg} driver. Tables
are Python classes:

\begin{lstlisting}[language=Python]
class JobAttempt(Base):
    __tablename__ = "job_attempts"
    id: Mapped[uuid.UUID] = mapped_column(Uuid, primary_key=True)
    job_id: Mapped[uuid.UUID] = mapped_column(ForeignKey("jobs.id"))
    fencing_token: Mapped[uuid.UUID] = mapped_column(Uuid, unique=True)
    lease_expires_at: Mapped[datetime] = mapped_column(DateTime(timezone=True))
    observed_state: Mapped[str] = mapped_column(String(32))
\end{lstlisting}

\textbf{Why async?} A long-poll claim spends most of its 25 seconds asleep.
With \texttt{async}/\texttt{await}, one process serves hundreds of waiting
agents, because a sleeping request costs almost nothing. With one thread per
request, sixty agents would need sixty threads doing nothing.

\paragraph{A trap we hit.} An async engine keeps a pool of connections that
belong to the event loop that created them. The operations tool first called
\texttt{asyncio.run()} twice in a row; the second call started a new loop and
reused pooled connections from the first, failing with ``attached to a
different loop''. The fix was to run both steps inside one
\texttt{asyncio.run()}.

\subsection{Schema migrations with Alembic}

Early on the app created tables at startup with \texttt{create\_all()} and
patched new columns in with hand-written \texttt{ALTER TABLE ... IF NOT
EXISTS}. That works until the first real deployment, then becomes dangerous:
there is no record of which changes a database already has.

\textbf{Alembic} keeps an ordered list of migration scripts and a one-row table,
\texttt{alembic\_version}, recording which one a database has reached:

\begin{verbatim}
control/migrations/versions/
  0001_baseline_schema.py      every table as it was
  0002_job_idempotency_key.py  adds jobs.idempotency_key + unique index
  0003_node_update_ring.py     adds nodes.update_ring
\end{verbatim}

\begin{lstlisting}[language=Python]
# 0003_node_update_ring.py
revision = "0003"
down_revision = "0002"

def upgrade() -> None:
    op.add_column("nodes", sa.Column("update_ring", sa.Integer(),
                  server_default=sa.text("2"), nullable=False))

def downgrade() -> None:
    op.drop_column("nodes", "update_ring")
\end{lstlisting}

Adding a column is now: edit \texttt{models.py}, run
\verb|alembic revision --autogenerate -m "..."|, review the generated
script, commit. The API applies pending migrations on startup.

\paragraph{Adopting migrations on an existing database.} Databases created the
old way have the tables but no \texttt{alembic\_version} row. On startup,
\texttt{init\_db} detects that case, applies the old column patches, and
\emph{stamps} the database at revision \texttt{0001} (records it as already
there) before upgrading. This was tested on a copy of the real development
database, then \texttt{alembic check} confirmed the schema matched the models
exactly.

\paragraph{A detail autogenerate got wrong.} The \texttt{jobs} and
\texttt{job\_attempts} tables reference each other. Autogenerate put both
foreign keys inline in \texttt{CREATE TABLE jobs}, but \texttt{job\_attempts}
did not exist yet. The baseline was edited to add those two keys at the end
with \texttt{op.create\_foreign\_key}. Always read generated migrations.

\paragraph{Other use cases.} Django migrations, Rails Active Record migrations,
Flyway and Liquibase in Java, Prisma Migrate in Node.

\section{The Control Plane API (FastAPI)}

\subsection{Routers and two kinds of caller}

The API has two audiences, split into separate router files:

\begin{description}
\item[Portal routes] (\texttt{routers/portal.py}) for people: submit jobs,
  list them, stream logs, admin actions. Authenticated by a session cookie from
  the university login, or the \texttt{X-Dev-User} header in development.
\item[Agent routes] (\texttt{routers/agent.py}) for machines: register,
  heartbeat, claim, renew, logs, outputs, complete. Authenticated by a
  per-node bearer token.
\end{description}

\subsection{Dependency injection}

FastAPI lets an endpoint declare what it needs as parameters, and builds those
values before calling it. Authentication is written once as a dependency:

\begin{lstlisting}[language=Python]
# control/app/deps.py (simplified)
async def get_current_node(
    session: AsyncSession = Depends(get_session),
    authorization: str | None = Header(default=None),
    x_boot_session_id: str | None = Header(default=None,
                                           alias="X-Boot-Session-Id"),
) -> Node:
    token = authorization.split(" ", 1)[1]
    cred = await session.scalar(select(NodeCredential).where(
        NodeCredential.token_hash == hash_token(token),
        NodeCredential.revoked_at.is_(None)))
    if cred is None:
        raise AppError(401, "auth", "invalid node token")
    node = await session.get(Node, cred.node_id)
    if x_boot_session_id != node.boot_session_id:
        raise AppError(409, "stale_boot_session", "re-register")
    return node

@router.post("/agent/heartbeat")
async def heartbeat(body: HeartbeatIn,
                    node: Node = Depends(get_current_node), ...):
    ...  # only runs for an authenticated, current node
\end{lstlisting}

Every agent endpoint gets authentication, revocation and boot-session checks
just by listing \texttt{node: Node = Depends(get\_current\_node)}.

\paragraph{Request validation.} Request bodies are Pydantic models
(\texttt{schemas.py}). A heartbeat with a missing field or a string where a
number belongs is rejected with \texttt{422} before any code runs.

\subsection{Long polling in code}

\begin{lstlisting}[language=Python]
@router.post("/agent/claim")
async def claim_job(node: Node = Depends(get_current_node),
                    session: AsyncSession = Depends(get_session)):
    waited = 0.0
    while True:
        attempt = await claim(session, node)
        if attempt:
            payload = await claim_payload(session, attempt)
            await session.commit()
            return {"attempt": payload, "build": None}
        await session.commit()          # release any locks while waiting
        if waited >= settings.claim_poll_seconds:
            return {"attempt": None, "build": None}
        await asyncio.sleep(1.0)        # costs almost nothing in async code
        waited += 1.0
\end{lstlisting}

The \texttt{commit()} before sleeping matters: holding a transaction open while
asleep would hold locks and a database connection for no reason.

\subsection{Live logs with Server-Sent Events}

The student's browser shows log lines as they arrive. The portal uses
\textbf{Server-Sent Events} (SSE): an ordinary HTTP response that never ends,
where the server writes one small text block per event:

\begin{verbatim}
id: 42
event: log
data: {"seq": 42, "stream": "stdout", "line": "epoch 3 loss 0.41"}

id: 43
event: log
data: {"seq": 43, "stream": "stdout", "line": "epoch 4 loss 0.37"}
\end{verbatim}

If the connection drops, the browser reconnects automatically and sends
\texttt{Last-Event-ID: 43}; the server resumes after line 43. That only works
because every log line has a sequence number, assigned by the agent and checked
by the server (a chunk must start at \texttt{last\_log\_seq + 1}).

\paragraph{SSE or WebSockets?} SSE is one-way (server to browser), plain HTTP,
and reconnects by itself, which is all a log viewer needs. WebSockets are
two-way, useful for chat or collaborative editing. Other SSE users: ChatGPT-style
token streaming, stock tickers, build logs in CI dashboards.

\subsection{Background workers}

Two loops start with the app (\texttt{main.py}):

\begin{itemize}
\item The \textbf{lease worker}, every 2 seconds: expire leases and mark silent
  nodes offline.
\item The \textbf{alert worker}, every 60 seconds: compute alerts and post new
  ones to a webhook.
\end{itemize}

They are \texttt{asyncio} tasks in the same process, created in FastAPI's
\texttt{lifespan} hook and cancelled on shutdown. For one server that is enough;
a larger system would move them to a separate worker process or a scheduler
like Celery beat.

\subsection{Errors as data}

All expected failures raise \texttt{AppError(status, code, message)}, which a
single handler turns into JSON like
\texttt{\{"error": "fenced", "message": "stale fencing token"\}}. Clients
branch on the stable \texttt{error} code, never on the human message. The agent,
for example, re-registers when it sees \texttt{stale\_boot\_session}.

\section{Scheduling: Which Job Goes Where}

\subsection{Matching a job to a node}

When a node claims, the scheduler walks the queue oldest first (higher priority
first) and takes the first job the node can run. ``Can run'' is
\texttt{fit\_failure()} in \texttt{scheduler.py}, which returns the first reason
a job does \emph{not} fit, or \texttt{None}:

\begin{lstlisting}[language=Python]
def fit_failure(node, job, env=None) -> str | None:
    res = job.resources
    if res.get("gpu_class") and res["gpu_class"] != node.compatibility_class:
        return "class"
    if node.vram_allocatable_mb < res.get("min_vram_mb", 0):
        return "vram"
    if res.get("cpu_cores") and node.cpu_cores < res["cpu_cores"]:
        return "cpu"
    if res.get("ram_mb") and node.ram_mb < res["ram_mb"]:
        return "ram"
    archs = [a for a in (env.gpu_archs if env else []) if a != "cpu"]
    if res.get("min_vram_mb") and archs and node_arch(node) not in archs:
        return "env"      # image not built/tested for this GPU
    if node.disk_free_bytes - watermark < res.get("disk_mb", 0) * 2**20:
        return "disk"
    return None
\end{lstlisting}

Returning a \emph{reason} instead of \texttt{True/False} pays off twice: the
claim loop can skip quickly, and the portal can tell the student \emph{why} a
job waits.

\subsection{Pending reasons people can read}

\texttt{explain\_pending()} runs when a student opens a job and produces one
sentence such as:

\begin{itemize}
\item \texttt{waiting for node class ada-12g} (no such PC exists or is online)
\item \texttt{waiting for a 12 GB GPU}
\item \texttt{waiting for a GPU this environment supports}
\item \texttt{waiting for disk on eligible nodes}
\item \texttt{waiting for user concurrency slot}
\item \texttt{waiting for a smaller GPU or high-VRAM reservation}
\end{itemize}

A queue that says ``pending'' with no reason generates support emails; a queue
that explains itself does not.

\subsection{Quotas}

Each user and each project has a maximum number of running jobs (defaults 2
and 4). Before taking a job the scheduler counts active attempts for that user
and project. This keeps one student with 50 queued jobs from occupying every
PC. Slurm calls this a QOS limit; Kubernetes calls it a ResourceQuota.

\subsection{Reserving big GPUs, and aging}

A 24 GB card is rare. If a small 4 GB job grabs it, a student needing 20 GB
waits for no good reason. The rule, in \texttt{reserved\_skip()}:

\begin{enumerate}
\item On a high-VRAM node, a small job is skipped while any large job that fits
  this node is waiting.
\item With no large job waiting, a small job may still use the big card after
  it has waited \texttt{aging\_steal\_seconds} (15 minutes). This is
  \textbf{aging}: anything that waits long enough eventually wins, so nothing
  starves forever.
\end{enumerate}

\begin{verbatim}
Queue: [8 GB job (age 2 min)]     Idle: 12 GB node, 24 GB node
  24 GB node claims -> skips (job is young, reserve the big card)
  12 GB node claims -> takes it                        (test_scheduling.py)
\end{verbatim}

\subsection{GPU classes}

Students do not know model names; they pick a class. The agent maps the GPU it
detects with \texttt{nvidia-smi} to a class using \texttt{agent/classes.json}:

\begin{lstlisting}[language=Python]
# agent/classify.py (simplified)
def classify_gpu(gpu_model: str | None, vram_mb: int) -> str:
    if not gpu_model or gpu_model.lower() == "none":
        return "cpu"
    arch = arch_for(gpu_model)        # "rtx 3050" -> "ampere"
    if not arch:
        return "unsupported"          # quarantined until added
    for row in catalog["classes"]:
        if row["arch"] == arch and row["min_vram_mb"] <= vram_mb <= row["max_vram_mb"]:
            return row["name"]        # e.g. "ampere-8g"
    return "unsupported"
\end{lstlisting}

The architecture matters as much as the memory: CUDA code is compiled for
specific \emph{compute capabilities} (8.9 for RTX 40xx ``Ada'', 8.6 for RTX 30xx
``Ampere'', 7.5 for GTX 16xx and RTX 20xx ``Turing''). Each environment lists
the capabilities it supports, and the scheduler only sends GPU jobs to nodes
whose capability is on that list.

\paragraph{A real case.} The first real Windows PC had a GTX 1660, which the
original catalog did not know, so it registered as \texttt{unsupported} and was
quarantined. Adding Turing meant three edits: new classes in
\texttt{classes.json}, compute capability 7.5 in the catalog images' list, and
an agent version bump. The PC then updated itself, re-registered as
\texttt{turing-6g}, and left quarantine without anyone touching it.

\paragraph{Other use cases.} Kubernetes node labels and \texttt{nodeSelector},
Slurm partitions and GRES (``generic resources'' like \texttt{gpu:a100:2}),
cloud instance types (\texttt{g5.xlarge}).

\section{Files: Object Storage with MinIO}

\subsection{What object storage is}

An object store keeps \emph{objects} (any bytes) under \emph{keys} (strings)
inside \emph{buckets}. There are no real folders, no editing in place, and no
file locking: you put a whole object, get it, or delete it. In exchange it is
simple, scales to billions of objects, and every major cloud speaks the same
protocol, Amazon S3's. MinIO is an open-source server that speaks S3, so the
same code works on a laptop and against AWS.

\begin{lstlisting}[language=Python]
from minio import Minio
client = Minio("127.0.0.1:9000", access_key="labminio",
               secret_key="labminio", secure=False)
client.put_object("lab", "jobs/7/promoted/model.pt", data, length)
client.get_object("lab", "jobs/7/promoted/model.pt").read()
\end{lstlisting}

\subsection{Key layout}

\begin{verbatim}
submissions/<job>/input.tar.gz                  what the student uploaded
jobs/<job>/attempts/<attempt>/logs/...          log chunks, per attempt
jobs/<job>/attempts/<attempt>/outputs/<path>    uploads from one attempt
jobs/<job>/promoted/<path>                      the accepted outputs
builds/<build>/log.txt                          pip output of an env build
\end{verbatim}

Keys are designed so that cleanup is a prefix delete and so that two attempts
can never overwrite each other.

\subsection{Promote only after fencing}

Outputs are uploaded under the \emph{attempt's} prefix. Only when
\texttt{/complete} passes the fencing check does the server copy them to
\texttt{promoted/}, after recomputing each file's SHA-256 and comparing it with
the checksum the agent reported. A zombie attempt may upload whatever it likes;
it lands under its own prefix and is never promoted. The same idea, ``write
somewhere private, then publish atomically'', appears in database write-ahead
logs, \texttt{git commit}, and blue/green deployments.

\subsection{Limits}

Each attempt has caps on log lines per second, total log bytes, file size and
total output size. Excess log lines are dropped with a single warning line
rather than failing the job; an oversized output fails the attempt with
category \texttt{resource} and is not promoted. Limits are how a shared
system survives one buggy script that prints a million lines a second.

\subsection{Garbage collection}

Nothing reads a failed attempt's outputs again.
\verb|python -m app.ops gc --days 30| deletes the \texttt{attempts/} prefixes of attempts that were never
accepted, and expires custom environment images no job has used recently. A
test proves it deletes a lost attempt's partial file while keeping the
promoted result.

\section{Containers, Images and the Registry}

\subsection{Why every job runs in a container}

A student's script is untrusted code on a shared PC. Running it directly
would let it read other students' files, fill the disk, or install software
that breaks the next job. A Docker container gives each attempt its own
filesystem view, process tree and resource limits, and the whole thing is
deleted afterwards.

\subsection{The hardened docker run}

\texttt{build\_docker\_run\_cmd()} in \texttt{agent/executor.py} builds this
command (abridged):

\begin{verbatim}
docker run --name lab-<attempt>
  --read-only                     root filesystem cannot be modified
  --network none                  no internet, no LAN, no exfiltration
  --pids-limit 512                fork bombs stop at 512 processes
  --cpus 1 --memory 512m          cgroup limits from the job's request
  --tmpfs /tmp:size=2g            scratch space in RAM, size-capped
  -v <input>:/input:ro            the student's code, read-only
  -v <output>:/output             the only writable place that is kept
  --user 10000:10000              not root (CPU jobs)
  --gpus all                      only when the job asked for VRAM
  <image> python /input/train.py
\end{verbatim}

The function also refuses to build any command that mounts
\texttt{docker.sock}: a container with the Docker socket can start a privileged
container and take over the PC. Each flag closes one door; together they make
a script able to compute and write results, and little else.

\paragraph{Cgroups.} \verb|--cpus| and \verb|--memory| are enforced by
Linux \emph{control groups}. A container over its memory limit is killed by the
kernel (exit code 137), which the agent reports as category \texttt{resource}.

\paragraph{GPUs in containers.} The NVIDIA Container Toolkit teaches Docker the
\verb|--gpus| flag: it mounts the GPU device and the host driver's libraries
into the container. Inside WSL2 the Windows driver is shared with Linux, which
is why lab PCs install the NVIDIA driver on Windows only.

\subsection{Tags and digests}

An image name like \texttt{pytorch/pytorch:2.5.1-cuda12.4-cudnn9-runtime} has
a \emph{tag}, which is a movable label: the publisher can point it at a
different image tomorrow. A \emph{digest} like
\texttt{registry:5000/lab-env/1bf18b82@sha256:2f7e3116...} is the SHA-256 of the
image manifest. Change one byte of the image and the digest changes, so pulling
by digest always gives exactly the same bytes.

The platform insists on digests for built environments: a build only counts
as ready when the builder reports a \texttt{repo/name@sha256:...} reference from
a registry. Two lab PCs running the same job therefore run the same software,
and a rerun next month runs what ran today. This is \emph{content addressing},
the same idea behind Git commit hashes and IPFS.

\subsection{Custom environments}

A student who needs \texttt{transformers==4.44.0} submits a
\texttt{requirements.txt}. Building it on every PC would be slow and could give
different results, so it is built \emph{once}:

\begin{enumerate}
\item The control plane normalises the requirements (sorted, lower-cased,
  duplicates removed; no \texttt{-r} or \verb|--index-url|) and hashes them
  together with the base image and GPU class:
\begin{lstlisting}[language=Python]
payload = {"requirements": reqs, "base": base_image, "compat": compat_class}
env_hash = sha256(json.dumps(payload, sort_keys=True).encode()).hexdigest()
\end{lstlisting}
\item The same hash means the same environment. A second student with the same
  requirements reuses the first student's image.
\item A \emph{builder} node (\verb|--role build|) writes a Dockerfile
  (\texttt{FROM base}, \texttt{RUN pip install -r requirements.txt}), builds it,
  runs an import test (and a CUDA test for GPU classes), records
  \texttt{pip freeze} as the lockfile, and pushes to the registry.
\item Jobs wait in \texttt{building} until the push succeeds, then every node
  pulls the image by digest.
\end{enumerate}

A broken \texttt{requirements.txt} fails on the builder in seconds, with the pip
error in the build log, instead of failing on a GPU node after a long wait.

\subsection{A bug: images tied to an old address}

The hash ignores the registry's address. When the laptop running the registry
changed Wi-Fi networks, its IP changed, but old builds still said ``ready'' with
digests pointing at the old IP. New jobs reused them and every pull failed. The
fix: when a ready build's digest names a different registry host than the
current \texttt{REGISTRY\_URL}, the build is reset to \texttt{pending} and
rebuilt, both at submit time and when a node claims. The lesson: when you cache
by content hash, ask what is \emph{not} in the hash.

\subsection{Docker-in-Docker for the simulated lab}

The simulated lab runs five ``PCs'' as containers on one laptop. Each runs its
\emph{own} Docker daemon inside (the \texttt{docker:dind} image, with
\verb|--privileged|). That matters: each simulated PC has its own image cache,
so a job on \texttt{sim-pc-3} really pulls from the registry like a separate
machine would, instead of sharing the laptop's cache.

\paragraph{Other use cases.} Container isolation is how GitHub Actions,
Google Cloud Run, AWS Lambda and Kubernetes run code; digests are how
Kubernetes and Helm pin exact images in production; build-once-pull-everywhere
is how every CI/CD pipeline ships software.

\section{The Agent on Each Lab PC}

\subsection{Standard library only}

\texttt{agent/} imports nothing outside Python's standard library: HTTP with
\texttt{urllib.request}, JSON with \texttt{json}, processes with
\texttt{subprocess}, threads with \texttt{threading}. Installing it is copying
files; updating it is copying newer files. On a fleet of sixty PCs, ``no
\texttt{pip install}'' removes a whole class of failures (no network to PyPI, a
mirror down, a dependency that changed).

\subsection{The main loop}

\begin{lstlisting}[language=Python]
def run_forever(self):
    self.register()
    while True:
        status, hb = self.heartbeat()
        update = hb.get("agent_update")
        if update and self.maybe_update(update):
            sys.exit(updater.EXIT_UPDATED)   # launcher restarts new code
        if not self.healthy or self.busy or hb["command"] in {"drain",
                                                              "quarantine"}:
            time.sleep(5)
            continue
        status, payload = self.claim()       # long poll
        if payload.get("attempt"):
            self.execute(payload["attempt"])
        elif payload.get("build"):
            self.execute_build(payload["build"])
\end{lstlisting}

While a job runs, a second thread renews the lease every few seconds. If a
renew response says \texttt{desired\_state: cancel} (the student pressed Cancel,
or the server's runtime limit passed), the agent sets a \texttt{stop} event and
the executor stops the container.

\subsection{A real bug: a blocking read}

The executor first read job output like this:

\begin{lstlisting}[language=Python]
# before
while True:
    if stop.is_set():
        kill_container(); return 143, "cancelled"
    line = proc.stdout.readline()      # waits here until a line arrives
    if line:
        on_line(line)
    elif proc.poll() is not None:
        break
    if time.time() > deadline:
        kill_container(); return 124, "timeout"
\end{lstlisting}

\texttt{readline()} \emph{blocks}: it does not return until the process prints
another line or exits. A training script that prints ``starting'' and then
trains silently for an hour leaves the loop stuck inside \texttt{readline()}
for that hour. During that time it never checks \texttt{stop} or the deadline,
so \textbf{Cancel did nothing and time limits were ignored}. The unit tests
used chatty scripts and never noticed; the simulated failure exam, which runs a
silent \texttt{time.sleep(120)} job, caught it.

The fix moves reading to a background thread that feeds a queue. The main loop
waits on the queue for at most half a second, so it checks \texttt{stop} and the
deadline twice a second no matter what the job prints:

\begin{lstlisting}[language=Python]
# after (agent/executor.py)
def _pump(proc, on_line, stop, deadline) -> str:
    lines = queue.Queue()

    def reader():
        for line in proc.stdout:
            lines.put(line.rstrip("\n"))
        lines.put(None)                 # end of output

    threading.Thread(target=reader, daemon=True).start()
    while True:
        if stop.is_set():
            return "stop"
        if time.time() > deadline:
            return "timeout"
        try:
            line = lines.get(timeout=0.5)
        except queue.Empty:
            continue
        if line is None:
            proc.wait()
            return "exited"
        on_line(line)
\end{lstlisting}

\paragraph{The general lesson.} Any loop that must react to more than one thing
(a cancel flag, a timer, incoming data) must never block on just one of them.
Other ways to solve the same problem: \texttt{select()} or \texttt{poll()} on
file descriptors, \texttt{asyncio} with timeouts, or Go's \texttt{select} over
channels. The regression tests start a script that prints once and sleeps for
60 seconds, then assert that a 1-second timeout and a cancel both return within
a few seconds. On the old code they hang.

\subsection{Knobs for simulation}

Three environment variables let a GPU-less container pretend to be a lab PC:

\begin{description}
\item[\texttt{LAB\_FAKE\_GPU=ampere-8g:8192:8.6}] report this class, VRAM and
  compute capability. Jobs still run on CPU because \verb|--gpus| is only
  added when \texttt{nvidia-smi} works.
\item[\texttt{LAB\_FAKE\_BUSY\_FILE=/sim/busy}] while the file exists, report
  the desktop GPU as busy.
\item[\texttt{LAB\_FAKE\_DISK\_FREE\_FILE=/sim/disk\_free}] report the number
  in the file as free disk space.
\end{description}

Test seams like these, small, explicit and off by default, let the failure
exam toggle conditions with \texttt{docker exec sim-pc-3 touch /sim/busy}
instead of really filling a disk.

\section{Self-Updating Agents}

Updating sixty PCs by hand is not realistic. The agents update themselves, and
the design has to answer three questions: how does a PC know the update is
genuine, how does it switch without a broken half-state, and what if the new
version does not work?

\subsection{The flow}

\begin{verbatim}
 server                                    lab PC
 ------                                    ------
 packs agent/ into a tar.gz,
 signs it with its private key
                     heartbeat  <--------  agent 0.7.0
 "agent_update: 0.7.1, sha256" -------->
                     GET /agent/bundle <-- (only between jobs)
 bundle + signature  --------------------> verify signature with pinned key
                                           unpack to releases/0.7.1
                                           write marker "pending"
                                           switch current -> releases/0.7.1
                                           exit 75
                                           launcher restarts agent 0.7.1
                     register   <--------  0.7.1 registers OK
                                           delete "pending"  (confirmed)
\end{verbatim}

\subsection{Signatures: proving the bundle is genuine}

Lab PCs may talk to the server over plain HTTP on campus Wi-Fi. Anyone who can
intercept that traffic could otherwise serve a fake ``update'' and run code on
every PC. So the server signs each bundle with an \textbf{RSA private key}, and
each agent checks the signature with the matching \textbf{public key}, which it
pinned when it was installed.

\paragraph{How RSA signing works, in miniature.} The key holder has a modulus
$n$, a public exponent $e$ (usually 65537) and a secret exponent $d$. To sign:
hash the data with SHA-256, pad it to a standard shape, treat it as a number
$m$, and compute $s = m^d \bmod n$. Anyone can verify by computing
$s^e \bmod n$ and checking that the result equals the padded hash. Only the
holder of $d$ can produce an $s$ that passes. Python's built-in three-argument
\texttt{pow} does modular exponentiation, so verification needs no crypto
library at all:

\begin{lstlisting}[language=Python]
# agent/updater.py
def verify_signature(data: bytes, signature_hex: str, public_key: dict) -> bool:
    n, e = int(public_key["n"], 16), int(public_key["e"])
    sig = bytes.fromhex(signature_hex)
    k = (n.bit_length() + 7) // 8                    # key size in bytes
    if len(sig) != k:
        return False
    em = pow(int.from_bytes(sig, "big"), e, n).to_bytes(k, "big")
    digest = SHA256_DIGEST_INFO + hashlib.sha256(data).digest()
    expected = b"\x00\x01" + b"\xff" * (k - len(digest) - 3) + b"\x00" + digest
    return hmac.compare_digest(em, expected)          # constant-time compare
\end{lstlisting}

The padding (\verb|00 01 FF FF ... FF 00| followed by an ASN.1 header that names
SHA-256) is the PKCS\#1 v1.5 standard, so the server signs with the ordinary
\texttt{cryptography} library and the agent verifies with ten lines of
arithmetic. \texttt{hmac.compare\_digest} compares in constant time so the
check does not leak information through timing.

\paragraph{Why not an HMAC with the node token?} That was the first idea. But
over plain HTTP the node token travels in every request header, so an attacker
who can intercept traffic could compute the HMAC too. A public-key signature has
no shared secret on the wire.

\paragraph{Trust on first use.} The installer fetches the public key from
\texttt{/agent/bundle-key} while an administrator is watching, and the agent
also pins the first key a server sends if it has none. After that the key never
changes silently. SSH host keys work the same way (``The authenticity of host
\dots{} can't be established. Continue?'').

\subsection{Atomic switching}

\begin{verbatim}
/opt/lab-agent/
  releases/0.7.0/   agent.py ...
  releases/0.7.1/   agent.py ...
  current -> releases/0.7.1        (a symbolic link)
  launcher.sh
  pending                          (only while 0.7.1 is unproven)
\end{verbatim}

The new release is unpacked next to the old one, byte-compiled to catch syntax
errors, and only then activated by replacing the \texttt{current} symlink:

\begin{lstlisting}[language=Python]
tmp = root / "current.new"
tmp.symlink_to(Path("releases") / version)
os.replace(tmp, root / "current")      # atomic rename on POSIX
\end{lstlisting}

\texttt{os.replace} is an atomic rename: anyone looking at \texttt{current}
sees the old target or the new one, never a missing or half-written link. This
is the same release layout Capistrano made popular for web deploys.

\subsection{Rollback}

Before switching, the updater writes a \texttt{pending} marker naming the new
and previous versions. The new agent deletes it right after a successful
registration. The launcher, a small bash loop that runs the agent, watches for
the case where the agent keeps exiting while the marker is still there:

\begin{lstlisting}[language=bash]
while true; do
  bash "$ROOT/current/run.sh"; code=$?
  [ "$code" -eq 75 ] && continue            # updated: restart at once
  if [ -f "$ROOT/pending" ]; then
    fails=$((fails + 1))
    if [ "$fails" -ge 3 ]; then             # new release cannot start
      switch_current "$prev"                # roll back
      echo "$bad" >> "$ROOT/rejected"       # never take it again
    fi
  fi
  sleep 5
done
\end{lstlisting}

The \texttt{rejected} list matters: without it, the server would keep offering
the broken version and the PC would loop update, crash, rollback forever.

\subsection{Rollout rings}

Every node has an \texttt{update\_ring}: 0 for canary, 1 for an early group, 2
for everyone. The server only offers an update to nodes whose ring is at most
\texttt{AGENT\_ROLLOUT\_RING}. A release goes to one canary PC first, then a few,
then all. \texttt{AGENT\_MIN\_VERSION} can additionally stop very old agents
from taking work until they catch up.

\paragraph{It worked on real hardware.} Bumping the agent from 0.7.0 to 0.7.1
(the Turing support change) on the server was enough: the real Windows PC
downloaded, verified and installed the update between jobs, re-registered with
its new GPU class and left quarantine. In the simulated lab, all five PCs did
the same in under a minute.

\paragraph{Other use cases.} Chrome and Firefox updates, Android A/B system
updates (two partitions, boot the new one, fall back if it fails), Kubernetes
rolling updates with readiness checks, and feature-flag percentage rollouts
all combine ``verify, switch atomically, watch, roll back''.

\section{Security}

\subsection{Who is calling? Authentication}

\begin{description}
\item[Students and admins] sign in through the university's identity provider
  using \textbf{OpenID Connect} (OIDC): the browser is sent to the campus login
  page, the user signs in there, and the identity provider sends back a signed
  token saying who they are. The platform never sees passwords. For
  development a mock identity provider (\texttt{oidc\_dev.py}) accepts any
  e-mail, and an \texttt{X-Dev-User} header skips login entirely.
\item[Lab PCs] receive a random 32-byte token when they first register. The
  server stores only a \emph{hash} of it:
\end{description}

\begin{lstlisting}[language=Python]
# control/app/crypto.py
def hash_token(token: str) -> str:
    return hashlib.sha256(f"{settings.token_pepper}:{token}".encode()).hexdigest()
\end{lstlisting}

If the database leaks, the attacker gets hashes, not usable tokens. The
\emph{pepper} is a server-side secret mixed into every hash, so the hashes are
useless without it too. Revoking a stolen PC's token
(\texttt{POST /nodes/<id>/revoke}) sets \texttt{revoked\_at}, and the
authentication dependency rejects it from then on.

\subsection{Refusing to run unsafely}

Development defaults are convenient and dangerous: the dev header lets anyone
be anyone. Rather than trusting a checklist, the API checks itself at startup
when \texttt{ENVIRONMENT=production}:

\begin{lstlisting}[language=Python]
def production_problems(self) -> list[str]:
    problems = []
    if self.dev_auth:
        problems.append("DEV_AUTH must be 0 (the X-Dev-User header ...)")
    if len(self.token_pepper) < 24:
        problems.append("TOKEN_PEPPER must be a random secret of 24+ chars")
    if self.minio_secret_key in {"labminio", ""}:
        problems.append("MINIO_SECRET_KEY is the dev default")
    ...
    return problems
\end{lstlisting}

With any problem present, startup fails with a list of exactly what to fix. A
misconfigured server never comes up half-safe. This was tested by starting the
production container with \texttt{DEV\_AUTH=1} and a short pepper; it refused
and named both.

\subsection{HTTPS with Caddy}

In production, Caddy sits in front of everything as a reverse proxy and
terminates TLS for three hostnames (portal, API, registry). Two certificate
options:

\begin{itemize}
\item \textbf{ACME} (Let's Encrypt): real, publicly trusted certificates,
  when the hostnames are reachable from the internet.
\item \textbf{Internal CA} (\texttt{tls internal}): Caddy runs its own
  certificate authority. Machines must be told to trust its root certificate
  once; the lab PC installer's \verb|--ca-cert| flag does that for both the
  agent and Docker.
\end{itemize}

\subsection{Read for everyone, write for one}

Every PC must pull images; only the builder should push them. Caddy enforces
that with HTTP methods alone:

\begin{verbatim}
{$REGISTRY_HOST} {
    @write not method GET HEAD
    basic_auth @write {
        {$REGISTRY_PUSH_USER} {$REGISTRY_PUSH_HASH}
    }
    reverse_proxy registry:5000
}
\end{verbatim}

Tested end to end: an anonymous pull returned 200, an anonymous push 401, and
the builder's push 202.

\subsection{Least privilege on the PC}

The agent's Linux user may start the Docker daemon at boot, and nothing else as
root. The installer adds exactly one sudo rule and verifies it with
\texttt{visudo}:

\begin{verbatim}
lab ALL=(root) NOPASSWD: /usr/sbin/service docker start
\end{verbatim}

A test in a fresh Ubuntu container confirmed \texttt{sudo service docker start}
works for that user and \texttt{sudo id} is refused.

\subsection{Thinking like an attacker}

A short threat model for this system:

\begin{center}
\begin{tabular}{ll}
\textbf{Threat} & \textbf{Defence} \\
Student script reads others' data & container, read-only input, no network \\
Script attacks the PC & non-root user, no docker.sock, cgroup limits \\
Fake agent update over Wi-Fi & RSA signature, pinned public key \\
Stolen PC replays its token & token revocation, hashed tokens \\
Student impersonates another & OIDC login, dev header off in production \\
Anyone pushes a malicious image & registry push requires credentials \\
Database leak & only peppered token hashes stored \\
\end{tabular}
\end{center}

Writing a table like this, even a rough one, is the first step of any security
review (the formal versions are called STRIDE or attack trees).

\section{Testing: From Unit Tests to Chaos}

The project tests at four levels, each catching things the others miss.

\subsection{Level 1: pytest against real services}

The 78 tests in \texttt{tests/} talk to a real Postgres and a real MinIO, not
mocks. Locking behaviour, unique indexes and transactions only exist in the
real database; a mocked database would happily ``pass'' the claim-locking bug.

\paragraph{Isolation.} \texttt{tests/conftest.py} drops and recreates a
\texttt{lab\_test} database and uses a \texttt{lab-test} bucket on every run,
so tests never touch development data:

\begin{lstlisting}[language=Python]
async def _recreate_test_db() -> None:
    conn = await asyncpg.connect(admin_url)
    await conn.execute(f'DROP DATABASE IF EXISTS "{TEST_DB}" WITH (FORCE)')
    await conn.execute(f'CREATE DATABASE "{TEST_DB}"')
    await conn.close()

asyncio.run(_recreate_test_db())
os.environ["DATABASE_URL"] = with_db(BASE_URL, TEST_DB)  # before app import
\end{lstlisting}

The environment variable is set before \texttt{app.config} is imported,
because the settings object reads the environment once, at import time.

\paragraph{Fake agents.} \texttt{harness/client.py} provides a
\texttt{FakeAgent} that speaks the agent protocol through FastAPI's test
client. A test reads like a story:

\begin{lstlisting}[language=Python]
def test_stale_token_rejected(client):
    portal, agent, job_id = isolated(client)
    attempt = agent.claim()
    agent.run_to_complete(attempt)
    bad = agent.complete(attempt, fencing_token=stale_token())
    assert bad["status_code"] == 409
\end{lstlisting}

Each test invents its own GPU class name (\texttt{iso-<random>}), so tests
cannot steal each other's jobs even though they share one database.

\subsection{Level 2: regression tests that fail on the old code}

When a bug is fixed, its test is first run against the \emph{old} code to
prove it catches the bug. In this project that was done for the claim-locking
bug (old code: the second node got nothing), the silent-job bug (old code: hung
past a 15-second limit), the registry-address bug (old code: the job was handed
an unpullable digest) and the Turing change. A test that passes on both old and
new code proves nothing.

\subsection{Level 3: the simulated lab and chaos testing}

Unit tests use fake agents. The simulated lab uses \emph{real} agents in
\emph{real} containers, then breaks things on purpose, an approach called chaos
engineering (Netflix's Chaos Monkey made it famous):

\begin{center}
\begin{tabular}{ll}
\textbf{Scenario} & \textbf{How it breaks things} \\
kill\_agent & \texttt{docker exec pc pkill -f agent.py} mid-job \\
reboot & \texttt{docker restart pc} mid-job \\
network\_drop & \texttt{docker network disconnect} for a full lease \\
cancel\_running & cancel a job while it trains \\
disk\_full & write a tiny number to the fake-disk file \\
desktop\_busy & touch the fake-busy file \\
max\_runtime & a 120 s job with a 5 s limit \\
custom\_env\_via\_registry & build on one PC, pull and run on another \\
self\_update & start a PC on agent 0.0.1 \\
\end{tabular}
\end{center}

Each scenario then asserts the outcome the specification demands, through the
API: the job completed exactly once with one promoted output, it was retried on
another node, the stale attempt was fenced, the cut-off PC rejoined. Short
timers (10 s leases) make the whole exam take about five minutes. It found the
silent-job bug on its first run, after 66 unit tests had passed.

\subsection{Level 4: real hardware}

Simulation cannot show Windows boot behaviour, real GPU drivers or a campus
network. The first real PC found two problems no simulation would: Windows
PowerShell 5.1's stderr handling (see the Windows chapter) and a GPU model missing
from the class list.

Once those were fixed, the first real GPU job ran on that PC: PyTorch inside the
container reported \texttt{cuda\_available True} on an NVIDIA GeForce GTX 1660,
and the output file was promoted. Along the way the agent restarted in the
middle of the job; the server failed that attempt as \texttt{node\_lost} and the
retry on the same PC succeeded. The recovery path the chaos exam rehearses had
just happened for real.

\subsection{Continuous integration}

GitHub Actions runs on every push:

\begin{itemize}
\item \texttt{control}: start Postgres and MinIO with Docker Compose, install
  dependencies, run \texttt{ruff}, run \texttt{alembic check} (are the
  migrations in sync with the models?), run pytest.
\item \texttt{portal}: \texttt{npm ci} and a production \texttt{next build}.
\item \texttt{sim-exam}: nightly and on demand, the full simulated failure exam.
\end{itemize}

CI caught nothing new in this project, which is its job: it makes sure what
passed on a laptop also passes on a clean machine, every time.

\paragraph{The testing pyramid.} Many fast unit tests at the bottom, fewer
slower integration tests in the middle, a few slow end-to-end tests at the top.
This project follows that shape: 78 tests in 16 seconds, 13 chaos scenarios in
5 minutes, one real PC.

\section{Running It in Production}

\subsection{Docker Compose stacks}

Docker Compose describes several containers, their settings, networks and
volumes in one YAML file. The project has three:

\begin{description}
\item[\texttt{docker-compose.yml}] development services (Postgres, MinIO,
  registry); the API usually runs on the host for fast reloads.
\item[\texttt{sim/docker-compose.sim.yml}] an \emph{overlay} adding the
  containerized control plane and five simulated PCs, enabled with
  \verb|--profile sim|. Compose merges overlay files on top of the base file.
\item[\texttt{deploy/docker-compose.prod.yml}] the production stack:
  Postgres, MinIO, registry, API, portal and Caddy. Only Caddy publishes ports.
\end{description}

\paragraph{Required variables.} In the production file every secret is
written as \verb|${POSTGRES_PASSWORD:?set POSTGRES_PASSWORD}|. The
\texttt{:?} syntax makes Compose refuse to start, naming the variable, when it
is unset or empty. Forgetting a secret is a startup error, not a silent default.

\paragraph{Health checks and ordering.} Services declare health checks (for
example \texttt{pg\_isready} for Postgres), and dependants wait with
\texttt{condition: service\_healthy}. The API never starts before its database
accepts connections.

\subsection{Backups}

\texttt{deploy/backup.sh} saves three things:
\begin{itemize}
\item Postgres, with \texttt{pg\_dump -Fc} (a compressed, restorable dump);
\item the MinIO bucket, with \texttt{mc mirror};
\item the agent signing key. Losing it would mean every PC rejects every future
  update until reinstalled.
\end{itemize}

Registry images are deliberately left out: they can be rebuilt from their
requirements. A backup is only real once it has been restored, so the restore
script was tested: back up the development data, restore into a scratch
database and bucket, and compare. 306 jobs, 99 nodes and 481 objects matched.

\subsection{Alerts}

\texttt{control/app/alerts.py} computes the conditions an administrator must
hear about without watching a dashboard: a PC offline for ten minutes, low
disk, many lost leases in an hour, one PC failing repeatedly, a job queued too
long. New alerts are posted to a chat webhook, each with a one-hour cooldown so
a persistent problem does not spam.

\subsection{Runbooks}

\texttt{docs/runbooks.md} has twelve step-by-step procedures: stand up the
server, add a PC, canary an update, drain, quarantine, replace WSL, recover a
full disk, rotate credentials, back up and restore, clean up storage, daily
health check, and moving the server. A runbook is written so that someone
tired at 2 a.m. can follow it, and each step says how to check it worked.

\subsection{Deployment ladder}

\begin{enumerate}
\item Simulated lab passes the failure exam.
\item Control plane on an always-on server with production settings.
\item One canary PC: reboot without login, CUDA job, self-update, rollback.
\item The failure exam by hand on five real PCs.
\item Staged rollout 1, 5, 15, 30, 60 PCs, recording success rate and lease
  losses at each step.
\end{enumerate}

Each step limits how much can break at once: the same reasoning as rollout
rings, applied to the whole system.

\section{Windows, WSL and the One-Command Setup}

\subsection{Why WSL}

Lab PCs run Windows, but containers and NVIDIA's container tooling are Linux
technologies. \textbf{WSL2} (Windows Subsystem for Linux) runs a real Linux
kernel in a lightweight virtual machine, and NVIDIA's Windows driver exposes the
GPU inside it. The agent and Docker Engine run in WSL Ubuntu; Windows only has
to start WSL at boot.

\subsection{Starting without anyone logging in}

A Windows scheduled task (\texttt{supervisor/install.ps1}) runs at startup
\emph{as the installing user} with the S4U logon type, meaning ``run as this
user even when they are not logged in, without storing their password''.
\texttt{SYSTEM} cannot start WSL2 on most PCs, which is why it is not used.
Inside WSL, systemd starts Docker, and the launcher starts the agent.

\subsection{The bootstrap script}

\texttt{supervisor/bootstrap.ps1} turns a blank PC into a lab node with one
command. The server serves it with its own address filled in:

\begin{verbatim}
powershell -ExecutionPolicy Bypass -Command
  "& ([scriptblock]::Create((irm http://server:8080/install/bootstrap.ps1)))"
\end{verbatim}

\texttt{irm} downloads the script as text, \texttt{[scriptblock]::Create} turns
the text into code, and \texttt{\&} runs it. It enables WSL, installs Ubuntu
without the interactive first-run prompt, creates a \texttt{lab} user, downloads
the install package, runs the Linux installer inside WSL, checks the GPU from a
container, installs the boot task and waits until the PC registers.

\subsection{Three lessons from the first real Windows PC}

\paragraph{1. Quoting across two operating systems.} A command typed in
PowerShell passes through PowerShell's parser, Windows' command-line rules, then
the Linux shell inside WSL. Each layer has its own quoting. The script avoids
the problem by never asking a shell to parse arguments:
\verb|wsl.exe -d Ubuntu -u root --exec bash /mnt/c/lab/wsl-prep.sh lab|.
\verb|--exec| runs the program directly with plain arguments, and anything
complicated goes into a script file written with Unix line endings.

\paragraph{2. Windows PowerShell 5.1 and stderr.} The first run stopped at the
GPU check with ``Unable to find image python:3.12.7-slim locally''. That message
is harmless: Docker prints it to stderr just before downloading. But in Windows
PowerShell 5.1, when a native program's output is redirected, each stderr line
becomes an \emph{error record}, and the script used
\texttt{\$ErrorActionPreference = "Stop"}, so the first stderr line aborted
everything. PowerShell 7 behaves differently, which is why the test container
did not catch it. The fix runs those commands with errors non-fatal and judges
them by exit code:

\begin{lstlisting}[language=powershell]
function Invoke-Native([scriptblock]$Block) {
    $saved = $ErrorActionPreference
    $ErrorActionPreference = "Continue"
    try { & $Block } finally { $ErrorActionPreference = $saved }
}
$out = Invoke-Native { & wsl.exe -d Ubuntu --exec docker run --rm `
                         --gpus all python:3.12.7-slim nvidia-smi -L 2>&1 }
if ($LASTEXITCODE -eq 0) { "GPU passthrough works" }
\end{lstlisting}

\paragraph{3. Networks that block GitHub.} \verb|wsl --install -d Ubuntu|
first downloads Microsoft's list of distributions from
\texttt{raw.githubusercontent.com}. On the lab PC's network that failed. The
bootstrap now falls back to \verb|--web-download|, then to importing
Canonical's official WSL image from \texttt{cloud-images.ubuntu.com}, or a file
brought on a USB stick:

\begin{verbatim}
wsl --import Ubuntu C:\lab\wsl\Ubuntu ubuntu-noble-wsl-amd64.rootfs.tar.gz
\end{verbatim}

\subsection{Line endings}

Windows uses CRLF (\verb|\r\n|) at the end of lines; Linux uses LF
(\verb|\n|). A bash script with CRLF fails with confusing errors like
\verb|$'\r': command not found|. Git for Windows converts line endings on
checkout by default, so a repository cloned on Windows could break every shell
script. A \texttt{.gitattributes} file fixes it for good:

\begin{verbatim}
*.sh  text eol=lf      # always LF, whatever the OS
*.cmd text eol=crlf    # batch files want CRLF
\end{verbatim}

The server also strips CR from scripts when it builds the install package,
belt and braces.

\section{Debugging Stories}

Every bug below was real, found while building this project. Each one carries
a lesson that applies far beyond it.

\begin{description}
\item[The image that disappeared.] One morning \texttt{docker compose up} failed:
  \texttt{minio/minio} could no longer be pulled from Docker Hub. CI had passed
  in August only because the image existed then.
  \emph{Lesson:} external dependencies can vanish; pin images by digest and
  have a fallback source.

\item[Old builds, new address.] Ready environment images pointed at a registry
  IP from a previous Wi-Fi network.
  \emph{Lesson:} when caching by a content hash, list what is not in the hash.

\item[The queue that looked empty.] Locking every queued job during a claim made
  concurrent claims see nothing.
  \emph{Lesson:} lock the smallest thing for the shortest time. Code that is
  ``correct'' can still be wrong under concurrency, and only a concurrent test
  shows it.

\item[Cancel that did nothing.] A blocking \texttt{readline()} meant a quiet job
  could not be cancelled or timed out.
  \emph{Lesson:} a loop that must react to several things must never block on
  one. Test with realistic inputs: real training jobs are quiet.

\item[One disk-full failure that was not.] The \texttt{disk\_full} chaos
  scenario failed with the wrong reason. The cause was the previous scenario's
  stuck job (the readline bug) leaving the node unable to heartbeat.
  \emph{Lesson:} when several tests fail, look for one root cause first.

\item[Two event loops.] The operations tool called \texttt{asyncio.run()} twice
  and reused pooled connections across loops.
  \emph{Lesson:} async resources belong to the loop that made them.

\item[A shadowed name.] In the same tool, a local variable \texttt{gc} (the
  argument parser) hid the function \texttt{gc()}, failing with
  ``\texttt{'ArgumentParser' object is not callable}''.
  \emph{Lesson:} name variables after what they are (\texttt{gc\_cmd}), and run
  the code path at least once in a test.

\item[A false alarm.] A launcher test seemed to leave the API running. It had
  simply been checked two seconds before its timeout fired.
  \emph{Lesson:} before calling something a bug, check the timeline.

\item[Every new dev user is an admin.] A test expected a student to get
  \texttt{403} from an admin endpoint, but the test settings made every dev user
  an admin.
  \emph{Lesson:} when a test fails, question the test's assumptions as hard as
  the code.

\item[PowerShell 5.1's stderr.] Harmless Docker output aborted the Windows
  setup script.
  \emph{Lesson:} the version you test on matters; PowerShell 7 and 5.1 differ.

\item[Blocked GitHub.] \verb|wsl --install| could not fetch its distribution
  list on a campus network.
  \emph{Lesson:} anything that downloads needs a second source.

\item[An unknown GPU.] A GTX 1660 was quarantined as ``cuda preflight failed''
  when the real reason was ``not in the class list''.
  \emph{Lesson:} error messages should say the real reason; a misleading message
  costs more time than the bug.
\end{description}

\section{Try It Yourself}

Reading builds understanding; breaking things builds intuition. Each exercise
takes ten to thirty minutes.

\begin{enumerate}
\item \textbf{Watch a lease expire.} Start \texttt{make dev}, submit a job with
  \texttt{python3 scripts/smoke.py}, and while it runs kill the local agent
  (\texttt{pkill -f agent.py}). Query the job every few seconds:
  \texttt{curl -H "X-Dev-User: alice@uni.edu" localhost:8080/jobs/<id>}. Watch
  the attempt go \texttt{running}, \texttt{unknown}, \texttt{failed
  (lease\_lost)}, and a second attempt start.

\item \textbf{Be a zombie.} In a Python shell, use \texttt{harness.client.FakeAgent}
  to claim a job, let its lease expire, let a second fake agent complete it,
  then call \texttt{complete} with the first agent's token. Read the
  \texttt{409} body.

\item \textbf{See SKIP LOCKED.} Open two \texttt{psql} sessions. In the first,
  \texttt{BEGIN; SELECT id FROM jobs ... FOR UPDATE SKIP LOCKED LIMIT 1;} and
  do not commit. In the second run the same query. Note which row each gets.
  Now remove \texttt{SKIP LOCKED} from the second and see it wait.

\item \textbf{Stream logs by hand.}
  \texttt{curl -N -H "X-Dev-User: alice@uni.edu"
  localhost:8080/jobs/<id>/logs/stream} shows raw SSE. Stop it, then rerun with
  \texttt{-H "Last-Event-ID: 3"}.

\item \textbf{Verify a signature by hand.} Fetch
  \texttt{/agent/bundle-key} and \texttt{/agent/bundle} (with a node token),
  then call \texttt{agent.updater.verify\_signature}. Flip one byte of the bundle
  and verify again.

\item \textbf{Break an update on purpose.} In the simulated lab, add
  \texttt{raise SystemExit(1)} at the top of \texttt{agent/agent.py}, bump the
  version, rebuild only \texttt{sim-control}, and watch the PCs update, fail
  three times, roll back, and list the version in
  \texttt{/opt/lab-agent/rejected}. Then revert.

\item \textbf{Add a scenario to the exam.} Write a \texttt{sim/exam.py}
  scenario that pauses a PC (\texttt{docker pause}) for longer than a lease and
  checks what happens to its job. What should happen when it resumes?

\item \textbf{Add a migration.} Add a nullable \texttt{notes} column to
  \texttt{nodes}, autogenerate a migration, read it, run \texttt{make test}, and
  confirm \texttt{alembic check} is clean.
\end{enumerate}

\subsection{Where these ideas show up next}

\begin{description}
\item[HTCondor] a decades-old system built for exactly this, scavenging idle
  desktops for computation, with policies that respect the desktop's owner.
\item[Slurm] the standard scheduler on university supercomputers; partitions,
  QOS limits and GRES map to this project's classes and quotas.
\item[Kubernetes] node labels, resource requests, leases for leader election,
  rolling updates and readiness checks.
\item[CI systems] GitHub Actions and GitLab runners pull jobs, run them in
  containers, stream logs and upload artifacts, the same loop as this agent.
\item[Render farms and BOINC] distributing independent work units to many
  machines and validating the results.
\item[Payment APIs] idempotency keys and exactly-once effects.
\end{description}

\subsection{Further reading}

\begin{itemize}
\item Martin Kleppmann, \emph{Designing Data-Intensive Applications}, chapters
  on replication, transactions and ``the trouble with distributed systems''.
\item Martin Kleppmann, ``How to do distributed locking'' (fencing tokens).
\item The PostgreSQL manual, ``Explicit Locking'' and the \texttt{SELECT}
  locking clauses.
\item The Twelve-Factor App (\texttt{12factor.net}).
\item Google, \emph{Site Reliability Engineering}, on alerting, runbooks and
  release engineering.
\end{itemize}

\section{Reference}

\subsection{Commands}

\begin{verbatim}
make dev                     local stack (API :8080, portal :3000)
make test                    pytest against lab_test
make lint                    ruff
make sim / sim-exam / sim-down   simulated lab
python3 scripts/smoke.py     submit a test job

cd control && ../control/.venv/bin/alembic revision --autogenerate -m "..."
cd control && ../control/.venv/bin/alembic upgrade head

python -m app.ops status     node states, queue, alerts   (in control container)
python -m app.ops gc --days 30 [--dry-run]

deploy/backup.sh /srv/backups
deploy/restore.sh /srv/backups/<timestamp>
\end{verbatim}

\subsection{Agent API (lab PCs)}

\begin{center}
\begin{tabular}{lll}
\textbf{Method} & \textbf{Path} & \textbf{Purpose} \\
POST & /agent/register & get node id, token, boot session \\
POST & /agent/heartbeat & health in, commands and updates out \\
POST & /agent/claim & long poll for an attempt or a build \\
POST & /attempts/\{id\}/renew & extend lease; learn about cancel \\
POST & /attempts/\{id\}/logs & numbered log chunks \\
PUT & /attempts/\{id\}/outputs/\{path\} & upload one output file \\
POST & /attempts/\{id\}/complete & finish, with fencing token \\
GET & /agent/bundle-key, /agent/bundle & signed self-update \\
GET & /install/bootstrap.ps1 & one-command Windows setup \\
\end{tabular}
\end{center}

\subsection{Portal and admin API}

\begin{center}
\begin{tabular}{lll}
\textbf{Method} & \textbf{Path} & \textbf{Purpose} \\
POST & /jobs & submit (multipart; optional Idempotency-Key) \\
GET & /jobs/\{id\} & status with a readable pending reason \\
GET & /jobs/\{id\}/logs/stream & live logs (SSE) \\
POST & /jobs/\{id\}/cancel, /retry & control a job \\
GET & /nodes, /admin/alerts & admin views \\
POST & /nodes/\{id\}/drain, quarantine, release & node state \\
POST & /nodes/\{id\}/ring, revoke & rollout ring, kill a token \\
\end{tabular}
\end{center}

\subsection{Important settings}

\begin{center}
\begin{tabular}{ll}
\textbf{Variable} & \textbf{Meaning} \\
ENVIRONMENT & \texttt{production} enables startup safety checks \\
LEASE\_TTL\_SECONDS & lease length (30) \\
DISCONNECT\_GRACE\_SECONDS & wait after expiry before failing (60) \\
CLAIM\_POLL\_SECONDS & long-poll length (25) \\
REGISTRY\_URL & registry host:port for custom environments \\
AGENT\_ROLLOUT\_RING & 0 canary, 1 early, 2 everyone \\
AGENT\_MIN\_VERSION & agents below this get no work \\
ALERT\_WEBHOOK\_URL & where alerts are posted \\
TOKEN\_PEPPER & secret mixed into token hashes \\
\end{tabular}
\end{center}

\subsection{Glossary}

\begin{description}
\item[ACME] protocol for obtaining TLS certificates automatically (Let's Encrypt).
\item[cgroup] Linux kernel feature that limits CPU, memory and processes.
\item[Compute capability] NVIDIA's version number for a GPU architecture.
\item[Digest] SHA-256 of an image manifest; an immutable image name.
\item[Fencing token] value that lets a server reject a stale owner.
\item[Idempotent] safe to repeat; repeating has no extra effect.
\item[Lease] ownership that expires unless renewed.
\item[Long poll] a request the server holds open until it has an answer.
\item[OIDC] OpenID Connect, login delegated to an identity provider.
\item[Pepper] server-side secret added to hashes.
\item[SSE] Server-Sent Events, a one-way stream over HTTP.
\item[TOFU] trust on first use; pin the first key seen.
\item[WSL2] Linux kernel running inside Windows.
\end{description}


\end{document}
